\chapter{Búsquedas}

\section{Búsqueda binaria en la respuesta}

\begin{lstlisting}[style=cppstyle, caption={Busqueda binaria en la respuesta}]
	bool can(int mid) {
		// Esta funcion debe responder si con mid como respuesta
		// es posible cumplir con la condicion del problema
		
		// Implementar la logica del chequeo
		
		return true;
	}
	
	int main(){
		// Busqueda binaria
		int low = 0, high = 1e9; // Ajustar limites segun el problema
		int ans = -1;
		
		while (low <= high) {
			int mid = (low + high) / 2;
			
			// si queremos el minimo que cumple
			if (check(mid)) {
				ans = mid;        
				high = mid - 1;
			} else {
				low = mid + 1;
			}
			
			// Si queres el maximo que cumple:
			if (check(mid)) {
				ans = mid;
				low = mid + 1;
			} else {
				high = mid - 1;
			}
		}
		
		cout << ans << endl;
	}
	
\end{lstlisting}
